\editionPage{donnees-utiles}{Données et relations utiles}{Toujours vérifier les unités et les conditions}{Valeurs arrondies pour les exercices, sauf indication contraire.}
\begin{tabularx}{\linewidth}{@{}YY@{}}\toprule
Constante ou grandeur & Valeur utilisée\\\midrule
Constante d'Avogadro & $N_A\simeq6,02\times10^{23}\unit{mol^{-1}}$\\
Charge élémentaire & $e\simeq1,60\times10^{-19}\unit{C}$\\
Masse d'un nucléon & $m_n\simeq1,67\times10^{-27}\unit{kg}$\\
Masse d'un électron & $m_e\simeq9,11\times10^{-31}\unit{kg}$\\
Constante de gravitation & $G\simeq6,67\times10^{-11}\unit{N.m^2.kg^{-2}}$\\
Pesanteur terrestre usuelle & $g\simeq9,8\unit{N.kg^{-1}}$ (ou 10 si fourni)\\
Lumière dans le vide ou l'air & $c\simeq3,00\times10^8\unit{m.s^{-1}}$\\
Son dans l'air ambiant & $v\simeq340\unit{m.s^{-1}}$\\
Eau liquide vers 20 °C & $\rho\simeq1,00\unit{g.mL^{-1}}$\\
Énergie massique de fusion de l'eau & $L_f\simeq334\unit{kJ.kg^{-1}}$\\\bottomrule
\end{tabularx}
\editionEtape{Conversions fréquentes}
$1\unit{L}=10^{-3}\unit{m^3}$ ; $1\unit{mL}=1\unit{cm^3}$ ; $1\unit{g}=10^{-3}\unit{kg}$ ; $1\unit{ms}=10^{-3}\unit{s}$ ; $1\unit{nm}=10^{-9}\unit{m}$. Pour passer de km/h à m/s, diviser par 3,6.
\editionEtape{Relations à retrouver avec leur sens}
$\rho=m/V$ ; $c_m=m_s/V_s$ ; $c_{m,m}V_p=c_{m,f}V_f$ ; $N=m/m_{\mathrm{entité}}$ ; $n=N/N_A$ ; $E=mL$ ; $P=mg$ ; $F=Gm_Am_B/d^2$ ; $v=d/\Delta t$ ; $f=1/T$ ; $n_1\sin i=n_2\sin r$ ; $U=RI$.
\editionRetenir{Une formule seule ne suffit pas : définir les symboles et vérifier le système, les hypothèses et les unités avant de l'appliquer.}
\editionPage{couverture-programme}{Correspondance avec le programme}{Repères de lecture et d'entraînement}{Le découpage du livre respecte les trois thèmes du BO et les compétences transversales.}
\editionEtape{01 · Mesurer et raisonner}
Unités et puissances de dix; Histogramme, moyenne, écart-type; Incertitude-type statistique; Chiffres significatifs et comparaison qualitative. Cours et exercices du chapitre 1, défi D01.
\editionEtape{02 · Identifier la matière}
Corps purs et mélanges; Tests eau, H2, O2, CO2; Masse volumique et températures; CCM; Composition massique et air. Cours et exercices du chapitre 2, défi D02.
\editionEtape{03 · Préparer et doser une solution}
Soluté, solvant, concentration maximale; Concentration en masse et masse volumique; Dissolution et dilution; Choix de verrerie; Étalonnage par teinte et masse volumique. Cours et exercices du chapitre 3, défi D03.
\editionEtape{04 · Atomes, noyaux et ions}
Atome, molécule, anion, cation; Composition du noyau et écriture A/Z; Masses et charges; Rapports d’échelles; Électroneutralité des composés ioniques. Cours et exercices du chapitre 4, défi D04.
\editionEtape{05 · Électrons et stabilité}
Configurations 1s, 2s, 2p, 3s, 3p; Position et familles du tableau périodique; Gaz nobles et ions usuels; Lewis fourni, doublets et stabilité; Énergie de liaison qualitative. Cours et exercices du chapitre 5, défi D05.
\editionEtape{06 · Compter les entités}
Masse d’une entité à partir de sa formule; Nombre d’entités N = m / m\_entité; Définition exacte de la mole; Quantité n = N / N\_A. Cours et exercices du chapitre 6, défi D06.
\editionEtape{07 · Changements d’état et énergie}
Six changements d’état et écriture; Modèle microscopique; Fusion et dissolution; Endothermique et exothermique; Énergie E = mL. Cours et exercices du chapitre 7, défi D07.
\editionEtape{08 · Réactions chimiques et bilans}
Réactifs, produits, spectateurs; Conservation des éléments et de la charge; Cinq familles de réactions usuelles; Stœchiométrie et réactif limitant; Effet thermique et masse du limitant. Cours et exercices du chapitre 8, défi D08.
\editionEtape{09 · Synthétiser et identifier}
Naturel et synthétique; Montage à reflux légendé; CCM légendée et comparaison; Isolement et contrôle; Sécurité et environnement. Cours et exercices du chapitre 9, défi D09.
\editionEtape{10 · Transformations nucléaires}
Isotopes; Écriture de réactions nucléaires; Fusion et fission; Soleil et centrales; Distinguer trois transformations. Cours et exercices du chapitre 10, défi D10.
\editionEtape{11 · Décrire un mouvement}
Système, point matériel et référentiel; Trajectoires et perte d’informations; Déplacement et vitesse moyenne vectorielle; Vitesse approchée et évolution; Python positions et vecteurs. Cours et exercices du chapitre 11, défi D11.
\editionEtape{12 · Forces et interactions}
Bilan des actions et vecteur force; Contact et distance; Gravitation vectorielle et poids; Fil et support; Actions réciproques. Cours et exercices du chapitre 12, défi D12.
\editionEtape{13 · Inertie et chute libre}
Principe d’inertie et conditions; Repos et mouvement rectiligne uniforme; Contraposée et variation de vitesse; Chute libre 1D avec ou sans vitesse initiale. Cours et exercices du chapitre 13, défi D13.
\editionEtape{14 · Émettre et transmettre un son}
Émission, vibration et résonance; Propagation et milieu matériel; Vitesse et mesure dans l’air; Retard, distance et aller-retour. Cours et exercices du chapitre 14, défi D14.
\editionEtape{15 · Période et perception sonore}
Signal périodique, période et fréquence; Domaines audibles, infrasons, ultrasons; Hauteur, timbre, amplitude; Échelle non linéaire des niveaux; Acquisition et production par microcontrôleur. Cours et exercices du chapitre 15, défi D15.
\editionEtape{16 · Lumière et spectres}
Propagation rectiligne et c; Longueur d’onde monochromatique; Spectre thermique continu; Spectres de raies et identification. Cours et exercices du chapitre 16, défi D16.
\editionEtape{17 · Réflexion et réfraction}
Réflexion et réfraction, plan d’incidence; Angles à la normale; Loi de Snell-Descartes; Détermination expérimentale d’un indice; Dispersion par prisme et réseau. Cours et exercices du chapitre 17, défi D17.
\editionEtape{18 · Lentilles et vision}
Lentille convergente, foyers et focale; Construction d’une image réelle; Grandissement et Thalès; Modèle de l’œil réduit; Formation expérimentale d’une image. Cours et exercices du chapitre 18, défi D18.
\editionEtape{19 · Circuits et lois électriques}
Mesurer tension et intensité; Lois des nœuds et des mailles, au plus deux mailles; Loi d’Ohm; Caractéristique et point de fonctionnement; Nuage de points et modélisation Python. Cours et exercices du chapitre 19, défi D19.
\editionEtape{20 · Capteurs et outils numériques}
Capteurs résistifs et exemples usuels; Courbe d’étalonnage d’une grandeur; Mesure avec un capteur; Microcontrôleur et chaîne d’acquisition; Python : nuage de points et modèle. Cours et exercices du chapitre 20, défi D20.

